\section{Related Literature}\label{sec:literature}

This paper sits at the intersection of three literatures: international standards coalitions, compatibility and interoperability in network competition, and policy-induced product differentiation. The contribution is not any one of these ingredients. It is the result-level consequence of allowing firms to reposition after a standards architecture is set: the induced product-market response can reverse governments' ranking of international standardization and a regional standards union relative to a fixed-position evaluation.

\paragraph{International standards coalitions and regional versus multilateral harmonization.}
The institutional benchmark is the literature on government standards policy and coalition formation. \citet{GandalShy2001} analyze a three-country environment with network effects, standardization unions, national welfare, and coalition incentives. \citet{TakaradaEtAl2020} study regional and multilateral standards harmonization with endogenous standards, national welfare, and core blocking, and show that a regional standards agreement can be the only core regime even when multilateral harmonization maximizes world welfare. \citet{KawabataTakarada2021} likewise examine deep trade agreements and standards harmonization. These papers already establish that regional standards coalitions, coalition-level harmonization, national-welfare participation incentives, and a national--global welfare wedge are economically relevant. None of those features is a setup contribution here. The additional margin in this paper is a separate horizontal product-positioning stage after standards policy, through which firms strategically re-differentiate and can change the government coalition ranking.

\paragraph{Compatibility, interoperability, and network competition.}
Continuous compatibility policy and firm-side differentiation are also established themes in industrial organization. \citet{Klimenko2009} studies technical-compatibility policy and international agreements under network externalities. \citet{BaakeBoom2001} analyze product differentiation, network externalities, and compatibility decisions; \citet{MatutesPadilla1994} show how partial network sharing can arise in a spatial banking environment; and \citet{DingKoShen2022} characterize partial compatibility and welfare in two-sided markets. More recent work studies interoperability directly. \citet{HuangEtAl2026} analyze weighted industry-wide and coalitional interoperability, pricing, profits, consumer surplus, and total welfare, while \citet{KretschmerEtAl2025} show that firms can strategically respond to mandated interoperability through data collection and privacy spillovers. This literature makes both the network-gain/competitive-insulation tradeoff and generic strategic response to interoperability regulation prior art. The mechanism studied here is narrower: a government standards architecture changes effective competitive distances and induces repositioning on a distinct horizontal product dimension, which then feeds back into national coalition incentives.

\paragraph{Policy-induced product differentiation.}
The closest predecessor for the timing from government standards policy to endogenous product characteristics is \citet{Ruiz2004}. Governments choose standard-recognition policy before firms choose product characteristics and prices, and endogenous differentiation can be strategically excessive. That result removes any novelty claim based simply on policy-induced firm adaptation or excessive differentiation. The distinction is instead in the equilibrium ranking. Ruiz's endogenous-characteristics extension preserves its qualitative government-policy result, whereas here the post-policy product-position response can reverse the member-country ranking between international standardization and a regional standards union relative to a fixed-position benchmark.

Taken together, the strongest hostile reconstruction combines \citet{GandalShy2001}, which supplies the three-country standards-union and national-welfare coalition architecture, \citet{Ruiz2004}, which supplies government standards policy followed by endogenous product characteristics, and \citet{TakaradaEtAl2020}, which supplies regional-versus-multilateral harmonization and blocking. That synthesis makes the setup unsurprising. The residual contribution is therefore deliberately stated at result level: holding the standards architecture fixed, allowing firms to reposition on a separate horizontal product dimension after policy can reverse member countries' $IS$--$SU$ welfare ranking relative to fixed product positions, making the regional union stable through producer-rent gains. At the canonical witness, fixing harmonization depths exogenously at their full-game values reproduces the reversal, so the paper does not claim that endogenous harmonization choice is necessary.

No priority claim is intended. The relevant statement is only that the audited closest literature does not directly establish this benchmarked post-policy repositioning reversal as a proposition or immediate corollary.